%% 1_intro
\section{Introduction}
\label{sec:intro}

Recovering a dense physical field from sparse, irregularly placed samples is a
recurring inverse problem in vision, and radio maps are an instance of it that is
unusually well posed for visual reasoning. A transmitter illuminates a scene, walls
absorb and diffract the wavefront, and the received strength at a point is largely
determined by what lies along the straight path back to the source. The structure is
the one that governs shading and cast shadows, and unlike most inverse problems the
geometry that produces it is directly observable in a building raster.

The field has responded to this problem the way vision usually does, by varying the
architecture. Convolutional encoder-decoders~\cite{levie2021radiounet},
transformers, state-space models, adversarial networks and diffusion models have all
been proposed for radio map estimation, each reporting an improvement over the last.
What has not been asked is whether architecture is the binding constraint. We ask
it directly, and the answer is that it is not.

Evaluating nine learned methods under one protocol on identical held-out scenes, we
find that they separate cleanly in line of sight and barely at all in non line of
sight. Line-of-sight error spans $5.18$ to $9.98$~dB across the panel, a coefficient of
variation of $20.1\%$, while non-line-of-sight error spans $11.12$ to $15.52$~dB at a
coefficient of variation of $9.6\%$. Discarding the weakest member tightens the
non-line-of-sight band to $11.12$ to $12.84$~dB at a coefficient of variation of $5.1\%$. Six
architecture families and two orders of magnitude of design effort separate these
methods, and in the regime that dominates the total error they are nearly
indistinguishable. A plateau that survives this much variation is not a capacity
limit. It is the signature of a quantity that none of these models can form from the
inputs it receives.

That quantity is a line integral. The fraction of the transmitter-to-pixel ray that
passes through building material is a soft shadow depth, it is fully determined by
inputs the networks already have, and it is nonetheless not something a stack of
local operators computes cheaply, because it is a long thin path-dependent sum
between two points that may be two hundred cells apart. Supplying it as one extra
input channel costs $432$ parameters and moves both error modes at once, taking
line-of-sight error to $3.85$~dB and non-line-of-sight error to $9.78$~dB, below the
band that every other method shares.

The consequences run against the usual intuition about scale. In a controlled
$2 \times 2$ design, geometry conditioning and a second network in series turn out to
be substitutes rather than complements. Measured separately they are worth $2.04$ and
$1.50$~dB, and applied together they deliver $2.87$~dB rather than the $3.54$~dB
their sum predicts. A plain single network given the geometry outperforms a cascade
denied it. Effort spent computing the right input buys more than effort spent adding
capacity.

The same finding explains a transfer failure that the benchmark literature has not
confronted. The line integral is anchored at the transmitter, and on real indoor
cellular measurements the transmitter is usually not locatable, because indoor
boosters share cell identifiers with the macro network. Across $67$ measured cells
only $17.9\%$ admit a reliable point source. We therefore tried to remove the anchor
by taking a minimum over measurement locations, and we report that it does not work.
Marginalising over anchors collapses the field toward zero and discards exactly the
directional structure that made it useful. On real measurements classical
interpolation consequently remains ahead of every learned reconstructor we trained,
at $5.51$ against $6.20$~dB with $p = 2.5 \times 10^{-30}$ over $1924$ held-out
points. The gap between simulated and deployed accuracy is not a domain shift that
fine-tuning closes. It is an input that one setting has and the other does not.

We validate the underlying premise on measurements rather than assuming it. At a
walled site, reference-point pairs separated by walls differ by an additional $2.31$
to $2.50$~dB at matched distance with disjoint confidence intervals, while at an
open car park used as a negative control the same measurement returns nothing and
cross-validation selects a wall penalty of exactly zero.

\noindent\textbf{Contributions.}
\begin{enumerate}[itemsep=1pt,topsep=3pt,parsep=0pt,leftmargin=1.2em]
    \item A diagnostic showing that nine learned radio map reconstructors spanning
    six architecture families share one non-line-of-sight plateau while differing
    widely in line of sight, which locates the bottleneck in the inputs rather than
    in the models.
    \item A ray-occlusion conditioning channel that removes the plateau for $432$
    parameters, and a $2 \times 2$ analysis establishing that it is a substitute for
    architectural depth rather than a complement to it.
    \item A negative result, with its mechanism identified, showing that the
    transmitter anchor cannot be marginalised away, together with the measurement
    evidence that the anchor is unavailable in four of five real cells.
    \item An evaluation on real indoor cellular crowdsensing with a negative control
    site, which quantifies how far the simulated and deployed regimes diverge and
    states explicitly where learning does and does not help.
\end{enumerate}
